\chapter{Liczby kardynalne Mahla i słaba zwartość}
\label{ch:mahlos-weak-compactness}
\chaptermark{Liczby kardynalne Mahla i słaba zwartość}

\section{Plan i założenia}
W Rozdziale~\ref{ch:large-cardinals} wykazaliśmy,
że mierzalność implikuje silną nieosiągalność,
oraz skonstruowaliśmy zanurzenie elementarne
$j:V\to M$ o punkcie krytycznym $\kappa$.
Tutaj wprowadzimy zbiory klubowe i stacjonarne,
liczby kardynalne Mahla, własność drzewa, słabą zwartość
oraz relacje podziałowe. Udowodnimy w ZFC, że
\[
 \kappa\ \text{jest mierzalną liczbą kardynalną}\quad\Longrightarrow\quad
 \begin{cases}
 \kappa\ \text{jest liczbą kardynalną Mahla},\\
 \kappa\ \text{jest słabo zwartą liczbą kardynalną},\\
 \kappa\longrightarrow(\kappa)^2_2.
 \end{cases}
\]
Nie postulujemy istnienia mierzalnych liczb kardynalnych.
Jeżeli nie wskazano inaczej, $\kappa$ oznacza
nieprzeliczalną regularną liczbę kardynalną, a liczby kardynalne
utożsamiamy z początkowymi liczbami porządkowymi.

\section{Zbiory klubowe i stacjonarne}

\begin{definicja}[Zbiór klubowy]\label{def:lc38-club}
Niech $\kappa$ będzie nieprzeliczalną regularną
liczbą kardynalną. Zbiór $C\subseteq\kappa$
jest \emph{nieograniczony}, jeżeli
$\forall\alpha<\kappa\ \exists\beta\in C$
takie, że $\alpha<\beta$.
Jest \emph{domknięty}, jeżeli dla każdej
niezerowej granicznej liczby porządkowej
$\delta<\kappa$ z nieograniczoności
$C\cap\delta$ w $\delta$ wynika $\delta\in C$.
Zbiór spełniający oba warunki jest
\emph{klubowy} (ang.\ \emph{club}).
\end{definicja}

\begin{definicja}[Zbiór stacjonarny]\label{def:lc38-stationary}
Dla nieprzeliczalnego regularnego $\kappa$
zbiór $S\subseteq\kappa$ jest
\emph{stacjonarny}, gdy
$S\cap C\ne\varnothing$ dla każdego
zbioru klubowego $C\subseteq\kappa$.
W przeciwnym razie jest
\emph{niestacjonarny}.
\end{definicja}

\begin{lemat}[Przecięcie mniej niż $\kappa$ zbiorów klubowych]
\label{lem:lc38-club-intersections}
Niech $\mu<\kappa$ i $(C_\xi)_{\xi<\mu}$
będzie rodziną zbiorów klubowych w
nieprzeliczalną regularną liczbą kardynalną $\kappa$.
Wtedy $\bigcap_{\xi<\mu}C_\xi$ jest klubowy.
\end{lemat}
\begin{proof}
Przecięcie jest domknięte, ponieważ
każdy jego punkt skupienia jest
punktem skupienia wszystkich $C_\xi$.
Dla $\mu=0$ teza jest oczywista.
Załóżmy $\mu>0$.
Ustalmy $\beta<\kappa$ oraz
$\beta_0>\beta$. Dla $\xi<\mu$
i $\gamma<\kappa$ określamy
$d_\xi(\gamma)=
 \min\{\delta\in C_\xi:\delta\ge\gamma\}$,
co jest możliwe dzięki
nieograniczoności $C_\xi$.
Rekurencyjnie przyjmujemy
\[
 \beta_{n+1}=
 \sup\bigl(\{\beta_n+1\}\cup
       \{d_\xi(\beta_n)+1:\xi<\mu\}\bigr).
\]
Regularność $\kappa$ zapewnia
$\beta_n<\kappa$ dla każdego $n<\omega$
oraz $\delta=\sup_{n<\omega}\beta_n<\kappa$.
Dla każdego $\xi<\mu$ mamy
$\beta_n\le d_\xi(\beta_n)<\beta_{n+1}$,
więc ciąg punktów $C_\xi$ o tych
wartościach jest współkońcowy w $\delta$.
Domkniętość $C_\xi$ daje
$\delta\in C_\xi$.
Zatem wspólne przecięcie
zawiera punkt $\delta>\beta$.
\end{proof}

\begin{definicja}[Przecięcie diagonalne]
\label{def:lc38-diagonal}
Dla rodziny $(C_\xi)_{\xi<\kappa}$
zbiorów w $\kappa$ określamy
\[
 \mathop{\triangle}_{\xi<\kappa}C_\xi
 =\{\alpha<\kappa:
       \forall\xi<\alpha\ (\alpha\in C_\xi)\}.
\]
\end{definicja}

\begin{lemat}[Diagonalne przecięcie zbiorów klubowych]
\label{lem:lc38-diagonal-club}
Jeżeli $(C_\xi)_{\xi<\kappa}$
jest rodziną zbiorów klubowych
w nieprzeliczalnej regularnej liczbie kardynalnej $\kappa$,
to $D=\mathop{\triangle}_{\xi<\kappa}C_\xi$
jest klubowy.
\end{lemat}
\begin{proof}
Niech graniczna $\delta<\kappa$
będzie punktem skupienia $D$.
Dla ustalonego $\xi<\delta$
wszystkie dostatecznie duże
$\alpha\in D\cap\delta$ należą do
$C_\xi$. Domkniętość $C_\xi$
daje $\delta\in C_\xi$.
Ponieważ $\xi$ było dowolne,
$\delta\in D$, co dowodzi domkniętości.

Ustalmy $\beta<\kappa$.
Konstruujemy rosnący ciąg
$\gamma_0>\beta$ oraz
\[
 \gamma_{n+1}\in
 \bigcap_{\xi\le\gamma_n}C_\xi,\qquad
 \gamma_{n+1}>\gamma_n.
\]
Przecięcie po prawej jest klubowe
na mocy Lematu~\ref{lem:lc38-club-intersections},
bo $\gamma_n+1<\kappa$.
Niech $\delta=\sup_n\gamma_n<\kappa$.
Dla każdego $\xi<\delta$
wszystkie dostatecznie późne
$\gamma_{n+1}$ należą do $C_\xi$;
ich supremum $\delta$ także należy
do $C_\xi$. Zatem $\delta\in D$
i $\delta>\beta$.
\end{proof}

\begin{definicja}[Funkcja regresywna]
\label{def:lc38-regressive}
Niech $S\subseteq\kappa$.
Funkcja $f:S\to\kappa$ jest
\emph{regresywna}, jeżeli
$f(\alpha)<\alpha$
dla wszystkich $\alpha\in S\setminus\{0\}$.
\end{definicja}

\begin{twierdzenie}[Lemat Fodora]\label{thm:lc38-fodor}
Niech $\kappa$ będzie nieprzeliczalną
regularną liczbą kardynalną, $S\subseteq\kappa$
zbiorem stacjonarnym, a $f:S\to\kappa$
funkcją regresywną.
Wtedy istnieje $\xi<\kappa$
takie, że $f^{-1}(\{\xi\})$
jest stacjonarny.
\end{twierdzenie}
\begin{proof}
Przecięcie zbioru stacjonarnego
z ogonem ponad $0$ pozostaje
stacjonarne, więc pomijamy argument $0$.
Przypuśćmy, że każde włókno
$S_\xi=f^{-1}(\{\xi\})$
jest niestacjonarne.
Dla każdego $\xi<\kappa$
wybieramy klubowy $C_\xi$
rozłączny z $S_\xi$.
ZFC pozwala wybrać je jednocześnie,
ponieważ indeksy tworzą zbiór.
Na mocy Lematu~\ref{lem:lc38-diagonal-club}
zbiór $D=\mathop{\triangle}_{\xi<\kappa}C_\xi$
jest klubowy.
Wybierzmy $\alpha\in S\cap D$,
$\alpha>0$, i połóżmy
$\xi=f(\alpha)<\alpha$.
Z definicji $D$ otrzymujemy
$\alpha\in C_\xi$, podczas gdy
$\alpha\in S_\xi$ wymaga
$\alpha\notin C_\xi$.
Jest to sprzeczność.
\end{proof}

\begin{przyklad}[Klub i zbiór niestacjonarny]
\label{ex:lc38-omega-one}
Dla $\kappa=\omega_1$ zbiór
niezerowych granicznych liczb
porządkowych poniżej $\omega_1$
jest klubowy: nad każdą
$\alpha<\omega_1$ leży
$\alpha+\omega<\omega_1$, a
granica granicznych liczb porządkowych
jest graniczna. Zbiór następników
$\{\alpha+1:\alpha<\omega_1\}$
jest z nim rozłączny,
a więc niestacjonarny.
W szczególności samo bycie
nieograniczonym nie oznacza
stacjonarności.
\end{przyklad}

\begin{uwaga}[Stacjonarność a normalność]
\label{rem:lc38-fodor-normal}
Lemat Fodora przypomina własność
normalnego ultrafiltru z
Rozdziału~\ref{ch:large-cardinals}.
Zbiory stacjonarne nie tworzą
jednak filtru: dwa takie zbiory
mogą być rozłączne.
Nie utożsamiamy stacjonarności
z przynależnością do jednej miary.
\end{uwaga}

\section{Liczby kardynalne Mahla}

\begin{definicja}[Liczba kardynalna Mahla]\label{def:lc38-mahlo}
Nieprzeliczalną liczbę kardynalną $\kappa$
nazywamy \emph{liczbą kardynalną Mahla},
gdy jest silnie nieosiągalna i zbiór
\[
 I_\kappa=\{\lambda<\kappa:
      \lambda\text{ jest silnie nieosiągalną liczbą kardynalną}\}
\]
jest stacjonarny w $\kappa$.
Nie wystarcza sama nieograniczoność
zbioru $I_\kappa$.
\end{definicja}

\begin{twierdzenie}[Liczba kardynalna Mahla jest granicą nieosiągalnych liczb kardynalnych]
\label{thm:lc38-mahlo-limit}
Jeżeli $\kappa$ jest liczbą kardynalną Mahla,
to dla każdego $\alpha<\kappa$
istnieje silnie nieosiągalna liczba kardynalna
$\lambda$ taki, że
$\alpha<\lambda<\kappa$.
Ponadto zbiór takich $\lambda$
przecina każdy zbiór klubowy w $\kappa$.
\end{twierdzenie}
\begin{proof}
Ogon $C_\alpha=\{\beta<\kappa:
\beta>\alpha\}$ jest klubowy:
jest nieograniczony i domknięty.
Z definicji stacjonarności
$I_\kappa\cap C_\alpha\ne\varnothing$.
Każdy element tego przecięcia
spełnia pierwszą część tezy.
Druga część jest dokładnie
warunkiem stacjonarności $I_\kappa$.
\end{proof}

\begin{uwaga}[Granica a stacjonarność]
\label{rem:lc38-mahlo-warning}
Własność „$\kappa$ jest granicą
nieosiągalnych liczb kardynalnych”
jest słabszym warunkiem niż
własność Mahla. Z faktu
nieograniczoności $I_\kappa$
nie wolno bez dowodu wyprowadzać
przecinania każdego klubu.
\end{uwaga}

\section{Własność drzewa i słaba zwartość}

\begin{definicja}[Drzewo, poziomy i gałąź]
\label{def:lc38-tree}
\emph{Drzewo} to zbiór częściowo
uporządkowany $(T,<_{T})$,
w którym ścisłe poprzedniki
każdego węzła $t$ są dobrze
uporządkowane przez $<_{T}$.
Ich typ porządkowy oznaczamy
$\operatorname{ht}_T(t)$.
Poziom $\alpha$ to
$T_\alpha=\{t\in T:
                 \operatorname{ht}_T(t)=\alpha\}$,
a wysokość drzewa to
$\sup\{\operatorname{ht}_T(t)+1:t\in T\}$.

Dla nieskończonej liczby kardynalnej
$\kappa$ \emph{drzewo $\kappa$-owe}
ma wysokość $\kappa$ i
$|T_\alpha|<\kappa$ dla
wszystkich $\alpha<\kappa$.
\emph{Gałąź współkońcowa} jest
łańcuchem przecinającym
każdy poziom $T_\alpha$.
\end{definicja}

\begin{definicja}[Słaba zwartość]
\label{def:lc38-weak-compact}
Nieprzeliczalna regularna
liczba kardynalna $\kappa$ ma
\emph{własność drzewa}, gdy każde
drzewo $\kappa$-owe ma
gałąź współkońcową.
Jest \emph{słabo zwarta},
gdy jest silnie nieosiągalna
i spełnia własność drzewa.
\end{definicja}

\begin{przyklad}[Drzewo binarne]
\label{ex:lc38-binary}
Niech $\kappa$ będzie silnie
nieosiągalny i
$T=2^{<\kappa}
 =\bigcup_{\alpha<\kappa}2^\alpha$,
uporządkowany przez właściwe
przedłużanie funkcji.
Każdy poziom ma moc
$2^{|\alpha|}<\kappa$,
więc jest to drzewo $\kappa$-owe.
Funkcje stale równe zeru
tworzą gałąź współkońcową.
Ten przykład nie rozstrzyga
własności drzewa: trzeba
uwzględniać wszystkie drzewa
$\kappa$-owe, nie tylko jedno.
\end{przyklad}

\begin{lemat}[Obrazy małych zbiorów]
\label{lem:lc38-small-image}
Niech $\kappa$ będzie mierzalną liczbą kardynalną,
a $j:V\to M$ zanurzeniem elementarnym
z Rozdziału~\ref{ch:large-cardinals}
o punkcie krytycznym $\kappa$.
Jeżeli $A$ jest zbiorem mocy
mniejszej od $\kappa$, to
$j(A)=\{j(a):a\in A\}$.
\end{lemat}
\begin{proof}
Dla $A=\varnothing$ jest to oczywiste.
W przeciwnym razie istnieje
$\mu<\kappa$ oraz surjekcja
$f:\mu\to A$. W $M$ funkcja
$j(f):j(\mu)\to j(A)$ jest
surjekcją. Ponieważ $\mu<\kappa$,
mamy $j(\mu)=\mu$.
Z elementarności
$j(f)(\xi)=j(f(\xi))$
dla każdego $\xi<\mu$.
Wszystkie elementy $j(A)$
są więc obrazami elementów $A$.
Odwrotna inkluzja jest
konsekwencją elementarności
relacji przynależności.
\end{proof}

\begin{twierdzenie}[Mierzalność implikuje własność drzewa]
\label{thm:lc38-measurable-tree}
Jeżeli liczba kardynalna $\kappa$ jest mierzalna,
to każde drzewo $\kappa$-owe
ma gałąź współkońcową.
W szczególności $\kappa$
jest słabo zwarty.
\end{twierdzenie}
\begin{proof}
Z Twierdzenia~\ref{thm:lc37-measurable-inaccessible}
wiemy, że $\kappa$ jest
silnie nieosiągalną liczbą kardynalną.
Niech $T$ będzie dowolnym
drzewem $\kappa$-owym i
$j:V\to M$ zanurzeniem
z punktem krytycznym $\kappa$
z Rozdziału~\ref{ch:large-cardinals}.
W $M$ drzewo $j(T)$ ma
wysokość $j(\kappa)>\kappa$.
Zatem na jego poziomie
$\kappa$ istnieje węzeł $t$.

Dla każdego $\alpha<\kappa$
węzeł $t$ ma dokładnie
jednego poprzednika $u_\alpha$
na poziomie $\alpha$.
Ponieważ $j(\alpha)=\alpha$,
elementarność daje
$(j(T))_\alpha=j(T_\alpha)$.
Z $|T_\alpha|<\kappa$
i Lematu~\ref{lem:lc38-small-image}
wynika
$j(T_\alpha)=\{j(s):s\in T_\alpha\}$.
Możemy więc jednoznacznie
wybrać $s_\alpha\in T_\alpha$
takie, że $j(s_\alpha)=u_\alpha$.

Dla $\alpha<\beta<\kappa$
mamy $u_\alpha<_{j(T)}u_\beta$,
a elementarność $j$ daje
$s_\alpha<_{T}s_\beta$.
Ponieważ $j$ jest definiowalny
z ustalonej miary, zastępowanie
w ZFC tworzy zbiór
$\{s_\alpha:\alpha<\kappa\}$.
Jest to gałąź przechodząca
przez wszystkie poziomy drzewa.
\end{proof}

\section{Miara normalna i stacjonarność nieosiągalnych}

\begin{lemat}[Kluby należą do miary normalnej]
\label{lem:lc38-clubs-in-measure}
Niech $\kappa$ będzie mierzalną liczbą kardynalną
i $\mathcal W$ niegłównym,
$\kappa$-zupełnym normalnym
ultrafiltrem na $\kappa$.
Każdy zbiór klubowy
$C\subseteq\kappa$ należy do
$\mathcal W$.
\end{lemat}
\begin{proof}
Gdyby $A=\kappa\setminus C$
należał do $\mathcal W$,
wybierzmy $\gamma_0\in C$
oraz $A'=A\cap\{\alpha:
\gamma_0<\alpha<\kappa\}$.
Z Lematu~\ref{lem:lc37-uniform}
ogon należy do $\mathcal W$,
więc $A'\in\mathcal W$.
Określmy
$f(\alpha)=\sup(C\cap\alpha)$,
$\alpha\in A'$.
Ponieważ $C$ jest domknięty,
a $\alpha\notin C$,
zachodzi $f(\alpha)<\alpha$.
Normalność daje
$B\in\mathcal W$, $B\subseteq A'$
i $\xi<\kappa$, dla których
$f|_B$ jest stale równa $\xi$.
Wybieramy $\gamma\in C$
z $\gamma>\xi$ i
$\alpha\in B$ z $\alpha>\gamma$:
taki $\alpha$ istnieje,
ponieważ $B$ ma moc $\kappa$
na mocy Lematu~\ref{lem:lc37-uniform}.
Wtedy $\gamma\in C\cap\alpha$,
czyli $f(\alpha)\ge\gamma>\xi$,
co jest sprzecznością.
\end{proof}

\begin{twierdzenie}[Mierzalna liczba kardynalna jest liczbą kardynalną Mahla]
\label{thm:lc38-measurable-mahlo}
Jeżeli liczba kardynalna $\kappa$ jest mierzalna,
to zbiór silnie nieosiągalnych
liczb kardynalnych mniejszych od $\kappa$
jest stacjonarny w $\kappa$.
\end{twierdzenie}
\begin{proof}
Niech $j:V\to M$ będzie
zanurzeniem elementarnym
o punkcie krytycznym $\kappa$
z Rozdziału~\ref{ch:large-cardinals}.
Przez
$I_\kappa=\{\alpha<\kappa:
\alpha\text{ jest silnie nieosiągalną liczbą kardynalną}\}$
oznaczamy odpowiedni zbiór.
W Rozdziale~\ref{ch:large-cardinals}
skonstruowano normalny ultrafiltr
$\mathcal W=\{X\subseteq\kappa:
                        \kappa\in j(X)\}$.

Liczba kardynalna $\kappa$ jest silnie
nieosiągalną liczbą kardynalną w $V$ na mocy
Twierdzenia~\ref{thm:lc37-measurable-inaccessible}.
Jest także silnie nieosiągalną liczbą kardynalną
w tranzytywnym $M$.
Istotnie, funkcja świadcząca
w $M$ o singularności $\kappa$
byłaby taką samą funkcją w $V$,
co przeczyłoby regularności.
Gdyby $M$ uważał, że $\kappa$
nie jest liczbą kardynalną, odpowiednia
bijekcja z mniejszą liczbą
porządkową także należałaby do $V$.
Jeśli $M$ uważałby, że
$2^\lambda\ge\kappa$ dla
pewnego $\lambda<\kappa$,
to wybór w $M$ dawałby iniekcję
$\kappa\to\mathcal P^M(\lambda)$.
W $V$ byłaby to iniekcja
$\kappa\to\mathcal P^V(\lambda)$,
sprzeczna z silną granicznością.

Z elementarności wynika, że
$j(I_\kappa)$ jest w $M$
zbiorem silnie nieosiągalnych
liczb kardynalnych poniżej $j(\kappa)$.
Ponieważ $\kappa<j(\kappa)$
i $M$ uważa $\kappa$ za
silnie nieosiągalną liczbę kardynalną,
mamy $\kappa\in j(I_\kappa)$.
Stąd $I_\kappa\in\mathcal W$.
Każdy klub także należy do
$\mathcal W$ na mocy
Lematu~\ref{lem:lc38-clubs-in-measure}.
Właściwy filtr nie zawiera
dwóch rozłącznych zbiorów.
Dlatego $I_\kappa$ przecina
każdy klub i jest stacjonarny.
\end{proof}

\section{Kolorowania par}

\begin{definicja}[Relacja podziałowa]
\label{def:lc38-partition}
Niech $\kappa,\lambda,\theta$
będą liczbami kardynalnymi, a $r\ge1$
liczbą naturalną.
Przez $[\kappa]^r$ rozumiemy
zbiór $r$-elementowych
podzbiorów $\kappa$.
Zapis
\[
 \kappa\longrightarrow(\lambda)^r_\theta
\]
oznacza, że dla każdego
$c:[\kappa]^r\to\theta$
istnieje $H\subseteq\kappa$
o mocy $\lambda$, na którym
$c$ jest stałe dla wszystkich
argumentów z $[H]^r$.
Zbiór $H$ nazywamy
\emph{jednorodnym}.
\end{definicja}

\begin{twierdzenie}[Mierzalność i jednorodność par]
\label{thm:lc38-measurable-partition}
Jeśli liczba kardynalna $\kappa$ jest mierzalna,
to $\kappa\longrightarrow(\kappa)^2_2$.
\end{twierdzenie}
\begin{proof}
Niech $c:[\kappa]^2\to\{0,1\}$
i niech $\mathcal U$ będzie
niegłównym $\kappa$-zupełnym
ultrafiltrem na $\kappa$.
Dla $\alpha<\kappa$ i $i<2$
definiujemy
\[
 B_{\alpha,i}
 =\{\beta<\kappa:\beta>\alpha,\
                   c(\{\alpha,\beta\})=i\}.
\]
Ogon ponad $\alpha$ należy
do $\mathcal U$, więc dokładnie
jeden z dwóch rozłącznych
zbiorów $B_{\alpha,0}$,
$B_{\alpha,1}$ należy do
$\mathcal U$. Oznaczmy
jego kolor przez $i_\alpha$.
Jeden ze zbiorów
$X_i=\{\alpha<\kappa:i_\alpha=i\}$
należy do $\mathcal U$.
Ustalmy taki $i$.

Rekurencyjnie zbudujemy rosnący
ciąg $(\alpha_\xi)_{\xi<\kappa}$
o wyrazach w $X_i$.
Po wybraniu $\alpha_\eta$
dla $\eta<\xi$, $\xi<\kappa$,
rozważamy przecięcie
\[
 X_i\cap\bigcap_{\eta<\xi}
                      B_{\alpha_\eta,i}.
\]
Jest ono w $\mathcal U$,
ponieważ $\xi<\kappa$.
Nie może być ograniczone w
$\kappa$, bo wszystkie zbiory
ograniczone mają moc mniejszą
od $\kappa$, a nie należą do
$\mathcal U$ na mocy
Lematu~\ref{lem:lc37-uniform}.
Wybieramy najmniejszy element
przecięcia większy od wszystkich
wcześniejszych $\alpha_\eta$.
Zbiór
$H=\{\alpha_\xi:\xi<\kappa\}$
ma moc $\kappa$ i dla
$\eta<\xi$ spełnia
$\alpha_\xi\in B_{\alpha_\eta,i}$.
Zatem każde dwie różne
liczby z $H$ otrzymują
ten sam kolor $i$.
\end{proof}

\begin{uwaga}[Właściwa hierarchia implikacji]
\label{rem:lc38-implications}
Wykazaliśmy bezpośrednio, że
mierzalność implikuje słabą
zwartość, własność Mahla i
powyższą relację podziałową.
Nie dowodziliśmy odwrotnych
implikacji ani pełnych
charakterystyk słabej zwartości
przez logikę infinitarną.
Tego rodzaju równoważności
wymagają osobnego dowodu,
a nie samego rysunku hierarchii.
\end{uwaga}

\begin{wniosek}[Dwa następstwa mierzalności]
\label{cor:lc38-hierarchy}
Jeśli istnieje mierzalna liczba kardynalna,
to istnieje liczba kardynalna Mahla
i słabo zwarta liczba kardynalna.
Dokładniej, każda mierzalna liczba kardynalna ma obie te własności.
\end{wniosek}
\begin{proof}
Pierwsza własność wynika
z Twierdzenia~\ref{thm:lc38-measurable-mahlo},
druga z
Twierdzenia~\ref{thm:lc38-measurable-tree}
oraz z nieosiągalności
mierzalnych liczb kardynalnych
w Rozdziale~\ref{ch:large-cardinals}.
\end{proof}

\begin{uwaga}[Nie mylić implikacji z niesprzecznością]
\label{rem:lc38-consistency}
Nieosiągalność implikuje
istnienie zbiorowego
modelu ZFC, jak pokazano
w Twierdzeniu~\ref{thm:lc37-rank-model}.
Zatem niesprzeczna ZFC nie
dowodzi istnienia liczb kardynalnych
Mahla, słabo zwartych ani
mierzalnych: przeczyłoby to
drugiemu twierdzeniu Gödla.
Sam ten argument nie dowodzi
względnej niesprzeczności
żadnego z tych aksjomatów.
\end{uwaga}

\begin{uwaga}[Zakres formalnej kontroli]
\label{rem:lc38-formal}
Kompilacja LaTeX i skończone
testy drzew i relacji podziałowych
nie są weryfikacją przez jądro Lean.
Status formalizacji wymaga
rzeczywistej kompilacji kodu Lean.
\end{uwaga}
